% Options for packages loaded elsewhere
\PassOptionsToPackage{unicode}{hyperref}
\PassOptionsToPackage{hyphens}{url}
\PassOptionsToPackage{dvipsnames,svgnames,x11names}{xcolor}
%
\documentclass[
  11pt,
]{article}
\usepackage{amsmath,amssymb}
\usepackage{iftex}
\ifPDFTeX
  \usepackage[T1]{fontenc}
  \usepackage[utf8]{inputenc}
  \usepackage{textcomp} % provide euro and other symbols
\else % if luatex or xetex
  \usepackage{unicode-math} % this also loads fontspec
  \defaultfontfeatures{Scale=MatchLowercase}
  \defaultfontfeatures[\rmfamily]{Ligatures=TeX,Scale=1}
\fi
\usepackage{lmodern}
\ifPDFTeX\else
  % xetex/luatex font selection
  \setmainfont[]{DejaVu Serif}
  \setmonofont[]{DejaVu Sans Mono}
\fi
% Use upquote if available, for straight quotes in verbatim environments
\IfFileExists{upquote.sty}{\usepackage{upquote}}{}
\IfFileExists{microtype.sty}{% use microtype if available
  \usepackage[]{microtype}
  \UseMicrotypeSet[protrusion]{basicmath} % disable protrusion for tt fonts
}{}
\makeatletter
\@ifundefined{KOMAClassName}{% if non-KOMA class
  \IfFileExists{parskip.sty}{%
    \usepackage{parskip}
  }{% else
    \setlength{\parindent}{0pt}
    \setlength{\parskip}{6pt plus 2pt minus 1pt}}
}{% if KOMA class
  \KOMAoptions{parskip=half}}
\makeatother
\usepackage{xcolor}
\usepackage[margin=1in]{geometry}
\setlength{\emergencystretch}{3em} % prevent overfull lines
\providecommand{\tightlist}{%
  \setlength{\itemsep}{0pt}\setlength{\parskip}{0pt}}
\setcounter{secnumdepth}{-\maxdimen} % remove section numbering
\ifLuaTeX
  \usepackage{selnolig}  % disable illegal ligatures
\fi
\usepackage{bookmark}
\IfFileExists{xurl.sty}{\usepackage{xurl}}{} % add URL line breaks if available
\urlstyle{same}
\hypersetup{
  pdftitle={Balanced Pairs in Diagonal-Gated Two-Arm Branches},
  pdfauthor={Aafi Khan Malek},
  colorlinks=true,
  linkcolor={Maroon},
  filecolor={Maroon},
  citecolor={Blue},
  urlcolor={Blue},
  pdfcreator={LaTeX via pandoc}}

\title{Balanced Pairs in Diagonal-Gated Two-Arm Branches}
\author{Aafi Khan Malek}
\date{November 2026}

\ifPDFTeX
\DeclareUnicodeCharacter{0393}{\ensuremath{\Gamma}}
\DeclareUnicodeCharacter{039B}{\ensuremath{\Lambda}}
\DeclareUnicodeCharacter{03B3}{\ensuremath{\gamma}}
\DeclareUnicodeCharacter{03B4}{\ensuremath{\delta}}
\DeclareUnicodeCharacter{03C1}{\ensuremath{\rho}}
\DeclareUnicodeCharacter{03C3}{\ensuremath{\sigma}}
\DeclareUnicodeCharacter{03C4}{\ensuremath{\tau}}
\DeclareUnicodeCharacter{03C9}{\ensuremath{\omega}}
\DeclareUnicodeCharacter{2084}{\ensuremath{{}_4}}
\DeclareUnicodeCharacter{2208}{\ensuremath{\in}}
\DeclareUnicodeCharacter{2209}{\ensuremath{\notin}}
\DeclareUnicodeCharacter{2212}{\ensuremath{-}}
\DeclareUnicodeCharacter{221A}{\ensuremath{\surd}}
\DeclareUnicodeCharacter{222A}{\ensuremath{\cup}}
\DeclareUnicodeCharacter{2260}{\ensuremath{\ne}}
\DeclareUnicodeCharacter{2282}{\ensuremath{\subset}}
\DeclareUnicodeCharacter{22C3}{\ensuremath{\bigcup}}
\fi
\begin{document}
\maketitle

Independent Researcher, Ahmedabad, India ·
\href{mailto:aafimalek2032@gmail.com}{\nolinkurl{aafimalek2032@gmail.com}}.\\
This research note was developed with AI assistance.

This is a \textbf{written structural application} of the Lean-checked
actual-ranking theorem
\texttt{WidthThree.jointComparisonFrequency\_of\_branchRankSum}\allowbreak{}. It
gives unbounded intertwined examples which the earlier opposite-arm
degree condition misses. The geometry-to-rank proof below is not
Lean-checked, and this is not a result for all width-three posets.

\subsection{Setting and precise
statement}\label{setting-and-precise-statement}

Let \(P\) be a finite poset. Take incomparable \(a,b\) such that each
strict predecessor of \(a\) is strictly below \(b\). Write \(U\) for
strict upset and \(S=\{b\}\cup(U(b)\setminus U(a))\). Assume
\(p=\Pr(a<b)\le 1/2\) in the uniform distribution on all linear
extensions of \(P\). A balanced pair has comparison probability in
\([1/3,2/3]\).

Suppose \(S\) is the ordinal sum of an \(h\)-point chain, \(h\ge2\),
from \(b\) to \(q\), and a tail partitioned into chains
\(x_0<\cdots<x_{r-1}\) and \(y_0<\cdots<y_{t-1}\). For nonnegative
integers \(d,e\), require the diagonal gates \[
 x_k<y_{k+d},\qquad y_k<x_{k+e}
\] whenever both endpoints of the corresponding gate exist. There is no
restriction on \(r,t\), they need not be equal, and additional cross
relations are permitted. Empty arms are allowed. If \(d=e=0\) and both
arms are nonempty, the gate hypotheses are inconsistent; no such vacuous
case is needed.

\textbf{Theorem A (rank pooling only).} If \(D=\max(d,e)\le 2h-2\),
\(P\) has a balanced pair.

In fact, if \(D\le2h-1\), the only possible no-balance case would have
each arm's first point high, i.e.~both \(\Pr(a<x_0),\Pr(a<y_0)>2/3\).
Below we also remove that case and extend the numerical range using the
separate written conditional XYZ argument.

\subsection{Exact dependence on the prior ranking
results}\label{exact-dependence-on-the-prior-ranking-results}

Write \(\rho_R(z)\) for the one-based rank of \(z\) in \(R\) restricted
to \(S\), and \(F(z)=\Pr(a<z)\), for \(z\in S\). All these \(z\) are
incomparable with \(a\). Each actual ambient extension is counted once;
induced orders are not assumed equally likely, and no independence or
product factorization of the cut and induced order is used.

The existing branch-order swap gives
\(\Pr(\text{all of }S\text{ before }a)\le p\), since cut counts are
nonincreasing within each fixed induced order. The existing
single-target pooling result
\texttt{WidthThree.pairCount\_le\_add\_of\_pooledBranchPositions}\allowbreak{} gives
\(F(z)\le F(u)+F(v)\) when \(\rho_R(z)\le\rho_R(u)+\rho_R(v)\) for every
actual ranking \(R\). This follows for probabilities by dividing counts
by the common positive number of extensions. The Lean-checked
\texttt{WidthThree.jointComparisonFrequency\_of\_branchRankSum}\allowbreak{} gives \[
 F(x)+F(y)\le 2(F(u)+F(v))
 \quad\text{if}\quad
 \rho_R(x)+\rho_R(y)\le 2(\rho_R(u)+\rho_R(v))
 \quad\text{for every }R.
\] Repeated anchors \(u=v\) are allowed in both. For direct replay, the
companion ranking package has exact statements in
\texttt{JointBranchOrderDraft.lean} and typed audit in
\texttt{JointBranchOrderAudit.lean}; the unformalized step here is
supplying the universal rank hypothesis from the branch structure and
threading the no-balance contradiction.

\subsection{Geometry-to-rank proof}\label{geometry-to-rank-proof}

Assume no pair is balanced. Thus \(p<1/3\), and each \(z\in S\) is low
if \(F(z)<1/3\) or high if \(F(z)>2/3\); these are the only
possibilities. Monotonicity of \(F\) along poset chains makes low sets
initial segments. The initial \(h\)-point chain is low by induction: its
fixed rank \(k\ge2\) is at most twice its preceding rank \(k-1\). If one
arm is empty the whole branch is a chain, so the same induction makes
its maximum low. Its comparison with \(a\) and the final cut then
partition all extensions, of total probability at most
\(F(\max S)+p<1\), impossible. This also handles both arms empty. From
now on \(r,t>0\).

Let \(i,j\) be the low-prefix lengths on \(X,Y\). Put \(m=\max(i,j)\),
\(s=\min(i,j)\), and \(\Delta=|i-j|=m-s\). A last-low anchor on an arm
is its last low or \(q\) when its low prefix is empty. Choose the
following pair \(u,v\) of low anchors, \textbf{fixed independently of
\(R\)}:

\begin{itemize}
\tightlist
\item
  If \(s\ge\Delta\), use the two arms' respective last-low anchors. The
  earlier anchor has rank at least \(h+s\). The later anchor has rank at
  least \(h+i+j\), since by its position all \(i+j\) lows of the two
  arms have appeared. This includes the case one anchor is \(q\), and if
  \(i=j=0\) both anchors are \(q\) of rank \(h\).
\item
  If \(\Delta>s\), repeat the last low of the longer low prefix; its
  rank is at least \(h+m\). This arm is nonempty by \(\Delta>0\).
\end{itemize}

Set \(M=\max(s,\Delta)\). In both cases the selected anchors satisfy, in
every order, \[
 A_R=\rho_R(u)+\rho_R(v)\ge 2h+i+j+M.                 \tag{1}
\] In the second case use \(2m=i+j+\Delta\). This choice is why one does
not lose slack when the two low prefixes are imbalanced.

Both arms have a first high. If all tail points are low, the two
comparisons of \(a\) against their chain maxima, together with the final
cut, cover all extensions and have total probability \(<1\). If exactly
one arm has a high, each point preceding its first high in any induced
order belongs to the common chain or one of the two low prefixes. Its
rank is at most \(h+i+j+1\le A_R\), so single-target pooling contradicts
high status. Write the two first highs as \(\alpha=x_i,\ \beta=y_j\).
They are incomparable: if \(\alpha<\beta\), then all highs on the other
arm succeed \(\alpha\), so again its rank is at most \(h+i+j+1\le A_R\).
The symmetric argument handles \(\beta<\alpha\).

Let \(L_X\) count the high \(X\)-points incomparable with \(\beta\), and
let \(L_Y\) count the high \(Y\)-points incomparable with \(\alpha\). In
every induced order the earlier high rank is at most \(h+i+j+1\). If the
later head is \(\beta\), its other possible earlier points are the at most
\(L_X\) high \(X\)-points incomparable with \(\beta\): none can lie
strictly below \(\beta\) in the poset, as that would imply
\(\alpha<\beta\). The symmetric bound uses \(L_Y\). Consequently, \[
 \rho_R(\alpha)+\rho_R(\beta)
 \le 2h+2(i+j)+2+\max(L_X,L_Y).
\] If \(\max(L_X,L_Y)\le2h+2M-2\), the right side is at most \(2A_R\) by
(1), uniformly for every actual \(R\). Applying exactly the verified
frequency endpoint displayed above makes the two high probabilities sum
to at most twice the sum of two low probabilities, a contradiction.
Therefore \[
 \max(L_X,L_Y)\ge2h+2M-1.                              \tag{2}
\]

Now use the gates (finite endpoints matter). By the \(X\) chain,
\(y_j<x_k\) for every existing \(k\ge j+e\): if such \(k\) exists, then
\(x_{j+e}\) exists and has the required gate from \(y_j\). Thus any
counted \(k\) lies in \(i\le k<j+e\). Similarly counted high
\(Y\)-indices lie in \(j\le k<i+d\). Both ranges are nonempty because
the first highs are incomparable. With \(\delta=i-j\) as a signed
integer, \[
 L_X\le e-\delta,\qquad L_Y\le d+\delta,
 \qquad \max(L_X,L_Y)\le D+\Delta.                    \tag{3}
\] Truncation at \(r,t\) only lowers these upper bounds. No exact
description of the cross comparabilities or restrictions on extra edges
are needed.

Combining (2) and (3) gives the necessary condition \[
 \boxed{D\ge 2h-1+2M-\Delta.}                         \tag{4}
\] Since \(2M-\Delta\ge0\), this proves Theorem A. The last term can
vanish only if \(\Delta=0\) and \(s=0\): indeed \(2M-\Delta\ge\Delta\)
and, if \(\Delta=0\), it equals \(2s\). This proves the assertion at
\(D=2h-1\).

\subsection{Optional widening with conditional XYZ (separate
dependency)}\label{optional-widening-with-conditional-xyz-separate-dependency}

Use the written
\href{https://github.com/Aafimalek/novel-discoveries/tree/6add80bc727aef16d920be7ce39d8312fed99a87/width-three-posets/verification/conditional-xyz-two-arm-slack}{conditional-XYZ
two-arm argument}: if \(q\) is low, \(q<x,y\), and
\(2\min(\rho_R(x),\rho_R(y))\le 3\rho_R(q)\) in every order, then
\(x,y\) cannot both be high. This conditional XYZ step is \textbf{not}
part of the checked frequency theorem and is not claimed as
Lean-checked.

Here the earlier high has rank at most \(h+i+j+1\) and \(\rho_R(q)=h\).
Therefore no balance entails \(2(i+j+1)>h\), or
\(2s+\Delta=i+j\ge K=\lfloor h/2\rfloor\). The nonnegative integer
\(T=2M-\Delta\) always satisfies \(3T\ge2s+\Delta\): the difference is
\(4(s-\Delta)\) if \(s\ge\Delta\), and \(2(\Delta-s)\) otherwise. Thus
(4) requires \(D\ge2h-1+\lceil K/3\rceil\). In particular:

\textbf{Theorem B (with written conditional XYZ).} The same conclusion
holds if \(D\le2h-2+\lceil\lfloor h/2\rfloor/3\rceil\).

\subsection{An actual longer intertwined width-three
family}\label{an-actual-longer-intertwined-width-three-family}

Take \(h=2, r=t=n\ge3, d=e=2\), put in both gate families and the arm
orders, and take transitive closure. Every generating edge advances the
index, and the first cross edge on a path advances it by two. Thus \[
 x_i<y_j\iff j\ge i+2,\qquad y_j<x_i\iff i\ge j+2.
\] The reverse implications are the gate followed by the other arm
order. Consequently cross-arm incomparability is exactly \(|i-j|\le1\).
The incomparability graph of the tail is connected (the diagonal links
and the adjacent links connect consecutive diagonal pairs), bipartite
and has a degree-three interior vertex for every \(n\ge3\). Its only
bipartition up to swapping is the two arms, so repartition does not
lower the maximum cross-arm incomparability degree. Also the common
ordinal chain prefix really ends at two: the tail has two incomparable
minimal points \(x_0,y_0\). Hence the earlier \(2h-2=2\) local-degree
criterion cannot cover this family by repartitioning. The earlier
crossed-gate criterion with \(x_0<y_1\), or its arm swap, does not apply
either. For comparison, \(n=1,2\) has degree at most two, hence is not
an example of strict improvement; those cases are covered by Theorem A
nonetheless.

No single-head rule from two initial low anchors can get started.
Placing \(y_0,y_1\) then \(x_0\) first in the tail is legal (validate
all generators, or take the three-element ideal and extend it); it gives
\(\rho_R(x_0)=5>4=2\rho_R(q)\). The symmetric order does this for
\(y_0\). But the earlier of these two heads is always at rank 3 and the
later has rank at most 5, so their sum is at most 8, just within joint
pooling. This is the genuine joint step at the empty-low-prefix corner.

Adjoining an isolated \(a\) realizes an ambient poset of width exactly
three with the stipulated exclusive branch: two chains cover \(S\), and
\(a,x_0,y_0\) form an antichain. Exactly one of the \(2+2n+1\) equally
possible insertions of \(a\), for each branch order, puts it before
\(b\); thus \(p=1/(2n+3)\le1/2\) in this explicit realization. Theorems
A and B apply also to other ambient posets with the stated good pair and
branch; they do not assume \(a\) isolated.

The companion bounded exact audit checks the integer differences, cross
closure, uniqueness of the arm bipartition via connectivity, and the
small \(n\) cases. Neither proof depends on this finite audit.

\end{document}
